\section{Stepwise Dilution Schedule}
\label{sec:dilution}

Working concentrations are prepared from the 0.5\,M stock (and from one another, in
series) using the volumetric dilution relationship:

\begin{calcbox}{\faVial\ Dilution Formula}
\[
C_1 V_1 = C_2 V_2 \qquad\Longrightarrow\qquad V_1 = \frac{C_2 V_2}{C_1}
\]
where $C_1, V_1$ describe the source (higher concentration) solution and $C_2, V_2$
describe the target (diluted) solution.
\end{calcbox}

\begin{tabularx}{\textwidth}{@{}cXXccX@{}}
\toprule
\textbf{Step} & \textbf{Target} & \textbf{Source Solution} & \textbf{Source Vol.} & \textbf{Diluent Vol.} & \textbf{Mixing Method} \\
\midrule
1 & 10\,mM in 500\,mL & 0.5\,M EDTA stock (this worksheet) & 10.0\,mL & 490.0\,mL DI water & Stir plate, 5 min \\
\rowcolor{labteal!5}
2 & 1\,mM in 250\,mL & 10\,mM working solution (Step 1) & 25.0\,mL & 225.0\,mL DI water & Vortex, 30 s \\
3 & 0.1\,mM in 100\,mL & 1\,mM intermediate (Step 2) & 10.0\,mL & 90.0\,mL DI water & Vortex, 30 s \\
\bottomrule
\end{tabularx}

\begin{notebox}{\faPen\ Dilution Notes}
\begin{itemize}[leftmargin=1.3em, itemsep=2pt]
  \item Prepare dilutions in the order listed --- each step draws from the solution made in
        the step immediately above it, not from the original stock.
  \item Use diluent (DI water) drawn from the same batch recorded in
        Section~\ref{sec:reagentspec} so that all working solutions share a traceable
        water lot.
  \item Label each bottle with its concentration, preparation date, and preparer's initials
        \emph{immediately} after mixing, before moving on to the next dilution, to prevent
        mix-ups between visually identical solutions.
  \item Discard any partially used working dilution bottle at the expiry assigned in
        Section~\ref{sec:signoff}; do not top up an existing bottle with fresh solution.
\end{itemize}
\end{notebox}
